\section{Background}

\subsection{Digital twin definitions and scope boundaries}
Digital twins are widely discussed across domains, but the literature still contains various different definitions and inconsistent scope boundaries. \citeauthor{inamdar2024_dt_microelectronics} describe the DT concept as spanning data acquisition, modeling, and analytics, and discuss how digital twins are frequently framed as continuously updated virtual representations, often emphasizing high-fidelity models and data-driven components for monitoring, diagnostics, prognostics, and optimization services~\cite{inamdar2024_dt_microelectronics}.
Their review also highlights that many proposed architectures converge around recurring building blocks such as data ingestion from the physical asset, a virtual model, analysis or simulation modules, and feedback mechanisms~\cite{inamdar2024_dt_microelectronics}.

Another perspective comes from digital twins in high-stakes domains such as precision medicine, where verification, validation, and uncertainty quantification (VVUQ) are central requirements. \citeauthor{sel2025_vvuq_precision_dt} emphasize that digital twins should be interpreted as decision-support and predictive systems that are dynamically updated with data from their physical counterparts, and that VVUQ is necessary for ensuring safety and efficacy when digital twins influence interventions~\cite{sel2025_vvuq_precision_dt}.
These works motivate two boundaries that matter for framework design:  a digital twin is not only a model, but an operational artifact that is updated from the physical system, and the requirements on trustworthiness and validation increase with the criticality of decisions enabled by the twin~\cite{inamdar2024_dt_microelectronics,sel2025_vvuq_precision_dt}.

\subsection{Digital twin reference architectures and implementation patterns}
\subsubsection{Typical digital twin reference architecture}
Modern digital twin (DT) systems share a common reference architecture made out of several layers and/or components. In a digital twin, data from physical assets is continuously ingested and transmitted via messaging or streaming to a virtual twin model. This virtual twin performs analysis or simulation and then feeds insights or control actions back to the physical world. This establishes a bidirectional data pipeline between physical and virtual assets. The data pipeline enables real-time monitoring, simulations, and closed-loop control~\cite{shao2020framework, jones2020characterising, soori2023digital}. Survey studies characterize such architectures in terms of maturity levels (from basic data integration up to autonomous operation). Moreover, the importance of maintaining traceability across the asset’s lifecycle is highlighted.

Most industrial digital twin platforms today are implemented using service-oriented, cloud-native principles. They are often built as collections of microservices that expose functionality through web APIs and communicate via event-driven messaging for telemetry and notifications~\cite{Fielding2000REST}. Hybrid cloud/edge deployments are common to meet latency and bandwidth constraints, with edge components bridging factory-floor protocols into the digital twin data pipeline~\cite{shao2020framework, soori2023digital}. On the modeling side, digital twins often integrate both physics-based models and analytics. Techniques like reduced-order modeling allow reuse of component models as fast, low-fidelity proxies~\cite{Kapteyn2020}, while domain-specific calibration methods ensure the twin reflects real process behavior~\cite{zheng2020quality, wang2021assembly}. Some digital twins support co-simulation scenarios for safe testing of changes~\cite{schluse2017experimentable}. Modern digital twin implementations are moving toward modular and loosely coupled designs, but achieving seamless composition of components across different users and lifecycle stages remains difficult~\cite{Liu2024DTMMReview, jones2020characterising}.

\subsection{Limitations and motivation for modularization}
\subsubsection{Why monolithic digital twins break down}
Early and current digital twin platforms have often been monolithic: treating a large asset or system as a single, integrated twin. This design has significant drawbacks for reuse and maintainability. \citeauthor{Gil2024} observe that “current platforms and architectures consider mostly monolithic digital twins and offer little support for coupling and checking the consistency of the coupling”~\cite{Gil2024}. In reality, integrating a new machine or process into a monolithic digital twin can require extensive redevelopment of the entire twin rather than a simple configuration change, and ensuring consistency between coupled components becomes difficult~\cite{Gil2024}. As a result, monolithic digital twins tend to be rigid and brittle since a change in one part of the production system can cause a redesign of the entire twin. The section below highlights key limitations.


\paragraph{Rigid, unitary design} Monolithic digital twin architectures bundle data, models, and services into one centralized application, which makes it difficult to extract or reuse sub-models. Developers must manually ensure consistency when connecting any new component, an error-prone process~\cite{Gil2024} due to the lacking native support for modular composition. This poor support for coupling different digital components limits integration of new assets.

\paragraph{Scalability challenges} As manufacturing systems or networks grow in scope, a single unified twin may struggle to scale. Performance and maintainability degrade when one digital twin is expected to represent an entire factory or product line. Specialized architectural concepts are required to manage such complexity. \citeauthor{Visser2024} propose a hierarchical “type digital aggregate” approach to achieve scalability in an automotive production network~\cite{Visser2024}. In general, monolithic twins often cannot be easily scaled up to larger systems without a major redesign.

\paragraph{Inflexibility to change} Rigid digital twin implementations are hard to adapt to evolving requirements or incremental system changes. Monolithic twins tend to accumulate excessive complexity that obscures the core behavior of the system over time~\cite{Visser2024}. By contrast, a library of reusable component models would allow engineers to replace or reconfigure parts of the twin as needed. \citeauthor{Kapteyn2020} demonstrate that using a library of modular models provides an “expressive, efficient, and intuitive” way to perform complex model updates for a digital twin~\cite{Kapteyn2020}. Monolithic designs often require redesigning or re-implementing large portions of the twin even for small changes in the physical system since they lack flexibility.



These issues make traditional digital twin solutions a poor fit for dynamic settings where production lines are frequently reconfigured and new equipment is added incrementally. The inability to easily swap in new components or extend the twin in a modular fashion leads to high integration effort and maintenance cost.

\subsubsection{Challenges motivating modularization}
Beyond the drawbacks of monolithic design, researchers and practitioners have noted a range of broader challenges in current digital twin implementations. Key recurring challenges include:


\paragraph{Interoperability and semantic integration} Diverse data sources, models, and vendor tools lead to interoperability gaps. Misaligned data formats and semantics across lifecycle stages make end-to-end integration difficult~\cite{jones2020characterising, kritzinger2018digital, Pretel2024SLR, Liu2024DTMMReview}. Even when standard interfaces exist, ensuring that information has the same meaning across design, production, and operation remains a major effort~\cite{Human2023}.



\paragraph{Synchronization and real-time control} Keeping the physical and virtual twin synchronized in real or near-real time is non-trivial. Issues of scheduling, network latency, and determinism can cause lags or inconsistencies between the twin and the physical asset. Data loads presenting themselves in bursts and tail-latency effects in distributed systems make timely updates unpredictable~\cite{tan2024dtsync, DeanBarroso2013TailAtScale}. Pure cloud-based twins often struggle to meet strict real-time constraints without special handling of edge computing and scheduling.

\paragraph{Model fidelity and drift} Ensuring the DT remains an accurate surrogate of the physical system over time is challenging. High-fidelity models are computationally expensive, while simpler models risk deviating from reality. Continuous calibration and validation are needed to keep models “fit for purpose” as conditions evolve~\cite{zheng2020quality, wang2021assembly, Liu2024DTMMReview}. During the DT lifecycle, model drift can occur if updates are not properly managed.

\paragraph{Verification and validation (V\&V)} Many current digital twins lack rigorous verification, validation, and experimentability provisions. Without systematic V\&V, it is hard to trust DT predictions or safely update the twin. Providing sandbox environments to test changes or certain scenarios is still an emerging practice~\cite{schluse2017experimentable}, this complicates change management.

\paragraph{Lifecycle integration and maintainability} Propagating changes through multiple enterprise systems (PLM, MES, ERP) and maintaining consistency between different versions of models, data, and services is labor-intensive. Configuration drift and version mismatches across sub-systems frequently introduce defects~\cite{vrabic2018digital, shao2020framework, soori2023digital}.

\paragraph{Security and governance} Expanding the digital twin across IT/OT boundaries increases the attack surface and raises concerns about data ownership and governance. Sharing DT data or control among partners must respect security policies, which is complex to manage without clear modular boundaries~\cite{shao2019cybersecurity}.

\paragraph{Scalability and operational overhead} Operating large numbers of twins or very complex twins poses scalability challenges. Monitoring and orchestrating many interdependent twin components demands robust observability and reliability practices. Data streaming, storage, and computation costs can grow quickly, and overall system complexity increases with each additional coupled component~\cite{Liu2024DTMMReview, soori2023digital, DeanBarroso2013TailAtScale}.


These challenges have moved researchers to advocate for more modular, composable, and decentralized digital twin architectures. In particular, adopting libraries of interchangeable model components~\cite{Kapteyn2020} or applying Multi-Agent System (MAS) design principles~\cite{Kalyani2024Survey, Pretel2024SLR, Lehmann2023, Siatras2024} may be promising methods to improve adaptability, reusability, and scalability of digital twins.

\subsection{Integration paradigms for modular digital twins}
This section examines common integration paradigms for modular digital twins and their implications for runtime adaptability, contrasting component-based, service-oriented, and agent-oriented approaches.

\subsubsection{Component-based and static coupling}
Component-based integration frameworks for digital twins largely stem from co-simulation tools and coupling libraries that link heterogeneous models with predefined interfaces. For example, the Functional Mock-up Interface (FMI) standard allows different simulation components to be packaged as Functional Mock-up Units and executed together. This enables tool-agnostic model exchange and co-simulation across domains \cite{blochwitz2012fmi20}. Such approaches excel at multi-model integration and time-synchronized execution. A complex system can be composed of reusable sub-models, each handled by its native simulator, while a master algorithm coordinates data exchange and timing. This yields reliable and high-performance integration. FMI-based co-simulations and libraries like MUSCLE3 can scale from single machines to HPC clusters, orchestrating message-passing between components with low overhead \cite{muscle3docs}. MUSCLE3,  uses a declarative YAML configuration to link models in a multiscale simulation, separating model logic from coupling logic and supporting features such as creating checkpoints and performance profiling for robust operation \cite{muscle3docs}.

However, these static coupling approaches fall short in adaptivity for modular digital twin integration. They typically require the interaction topology and interfaces to be defined in advance and then remain fixed during runtime. As a result, the digital twin cannot easily incorporate new components, alter couplings, or change its coordination logic on the fly. For example, MUSCLE3's configuration is static once the simulation starts, which limits runtime structural changes and prevents spawning or retiring simulation entities based on unforeseen conditions \cite{muscle3docs}. Similarly, FMI-based systems rely on predetermined component contracts. Integrating a new sensor model or switching out a subsystem usually requires stopping the system and redeploying with a revised model lineup. In short, while co-simulation frameworks ensure consistency and performance, they lack the dynamic reconfiguration and autonomous decision-making needed for adaptive digital twins.

This limitation motivates the MAS-based architecture. By using agents to encapsulate and orchestrate simulation components, the goal is to retain modular coupling benefits of frameworks like FMI and MUSCLE3 while enabling runtime adaptation. An agent-mediated twin can rearrange connections or introduce new agents in response to environmental changes, achieving flexibility that static coupling libraries alone do not provide.

\subsubsection{Service-oriented integration (SOA, microservices, workflows)}
Service-Oriented Architecture (SOA) is a paradigm in which software functionality is packaged into service encapsulated components accessible via well-defined interfaces. Services typically perform a discrete function  and are invoked on demand by clients. In distributed systems, services are composed loosely coupled, often through remote procedure calls or APIs of web services~\cite{singh2005service}. In the context of DT, many platforms follow a service-oriented approach: each major function (data ingestion, simulation, visualization, etc.) is deployed as an independent service accessible through an API. This design emphasizes modularity in deployment, scalability through replication of services, and clear interface contracts for integration. However, services are fundamentally passive components since they execute their functionality only when externally called and do not inherently pursue objectives beyond responding to requests.

In many digital twin platforms, this service orientation manifests as workflows and microservice architectures. In these designs, each functional element of a twin, including data ingestion, analytics, simulation, and user interaction, is implemented as an independent service, communicating via APIs, message queues, or event buses. This microservice-oriented strategy has been adopted in both industry platforms and research prototypes. Academic efforts similarly propose service-based frameworks. \citeauthor{borodulin2017dtcloud} outline a microservices and workflow architecture for a smart-factory digital twin, in which a workflow engine coordinates services for simulation, data processing, and control~\cite{borodulin2017dtcloud}. Such approaches solve the problem of interoperability and deployment at scale. Each component can be developed and updated independently, and standard web protocols or streaming platforms handle integration, making it trivial to plug the digital twin into cloud-to-edge environments~\cite{bellavista2024microservices_serverless_dt}. Moreover, workflows can explicitly define complex twin processes, which is useful for transparency and manageability in scenarios like industrial procedures or long-running simulations.

Despite these strengths, service-oriented and workflow-driven solutions often lack the adaptive autonomy required for highly modular digital twin integration. Coordination of microservices is usually predefined by orchestration logic, typically a static workflow or a set of rules in an orchestrator. Changes in the twin's behavior are hard to accommodate without manually reconfiguring the workflow or deploying new service instances. The integration is only as flexible as the predetermined orchestration. There is no built-in mechanism for the system to reorganize itself dynamically. Another limitation is that pure microservice architectures do not inherently provide intelligent decision-making about which services to invoke or when to restructure the process. This can make it difficult to achieve on-line adaptation or optimization beyond scaling the number of instances. Furthermore, heavy reliance on cloud orchestrators can introduce vendor lock-in and latency issues for real-time control \cite{marah2024_madtwin}, which is problematic for certain cyber-physical deployments.

\subsubsection{Agent-oriented integration (autonomy and coordination)}
These shortcomings motivate a shift toward more autonomous coordination. A MAS-based approach addresses this by deploying agents as coordinating entities. Instead of a static workflow engine, agents can negotiate, coordinate, and adjust their interactions at runtime. The proposed architecture can use the modular deployment benefits of microservices, while embedding decision logic and adaptability in agent behaviors rather than fixed orchestration scripts. This supports runtime reconfiguration such as forming new communication pathways or reallocating tasks in response to change.

An agent, is an autonomous software entity situated in an environment, capable of making independent decisions and interacting with others to achieve its objectives~\cite{singh2005service}. Agents are characterized by properties such as autonomy (agents can act without direct external invocation) and social ability (agents communicate and coordinate with other agents)~\cite{singh2005service}. Rather than being activated by external calls, an agent continuously senses its environment or incoming messages, reasons about its goals, and acts proactively. For example, in a DT context an agent representing a machine might actively monitor the machine’s state and decide to adjust a process or send a warning, without a central controller telling it to do so. Agents encapsulate their own internal state and logic, and they are inherently goal-driven: each agent strives to accomplish specific tasks or objectives it is assigned (or that it formulates for itself). This goal-oriented, proactive behavior is a key distinction from service components, which do not have internal drives or long-lived objectives.

The architectural difference between services and agents has important implications for integrating complex DT systems. A service-based digital twin relies on external orchestration to accomplish system-wide behavior: a separate workflow or controller must invoke the appropriate services in the right sequence to, say, simulate a production scenario or adjust operations based on twin insights. On the other hand, an agent-based digital twin can show emergent coordination and more adaptive system behavior. Because agents can initiate actions and negotiate with each other, a network of agents can collectively solve problems or adapt to changes without requiring all decisions to funnel through a central orchestrator. This means that in dynamic or unpredictable manufacturing environments, agent-oriented twins can respond more easily to disturbances. As \citeauthor{Lehmann2023} argue, MAS architectures (with autonomous, goal-seeking agents) are well-suited for “intelligent digital twins” that need to be active and adaptive in real time~\cite{Lehmann2023}. Empirical studies have shown that distributing decision-making to cooperative agents can improve performance. For example, in a factory case study, decentralized scheduling agents were able to detect bottlenecks early and reconfigure production schedules on the fly~\cite{Siatras2024}. This level of agility would be difficult to achieve with a static, service-orchestrated twin. Moreover moving from one monolithic twin to a society of “autonomous digital units” has been identified as a key step toward more interoperable and flexible digital twin systems~\cite{David2023}. 

Services provide reliable, predictable building blocks invoked by external logic, whereas agents bring initiative and self-organization. For next-generation digital twins that must integrate many components and continuously adapt, the agent-oriented approach offers distinct advantages in flexibility and resilience. However, it does come with the cost of added complexity in ensuring global coherence.

\subsection{Agents and digital twins: separation of responsibilities}
Work at the intersection of digital twins and MAS has increasingly focused on conceptual clarity: what belongs inside the digital twin abstraction, and what is better handled by agents. \citeauthor{mariani2022_dt_agents_cross_fertilisation} discuss digital twins and agents as complementary abstractions, where the digital twin can be understood as a shared medium that provides access to and interaction with physical assets, while agents provide distributed reasoning and decision-making capabilities \cite{mariani2022_dt_agents_cross_fertilisation}.
A core contribution of this perspective is the separation of concerns between the digital twin as an artifact closely tied to the local cyber-physical context of a specific physical asset, and agents operating in a broader application context that can combine information and actions across multiple twins and external services~\cite{mariani2022_dt_agents_cross_fertilisation}.
This motivates architectures where digital twins encapsulate asset-facing interfaces and state synchronization, while agents implement control, optimization, coordination, and cross-asset reasoning on top of those interfaces~\cite{mariani2022_dt_agents_cross_fertilisation}.

\subsection{MAS Fundamentals for digital twin architecture}
Having contrasted service-oriented and agent-oriented integration at a high level, we now outline fundamental concepts from Multi-Agent Systems (MAS) that inform a modular digital twin architecture. In a MAS-based approach, each asset or functional unit in the physical system can be represented by its own \emph{agent} – essentially a mini digital twin with autonomous behavior – and these agents collectively model and control the overall system. We review core agent concepts, organizational structures, coordination mechanisms, and classifications of agent types, highlighting how these elements can be applied in a digital twin context.

\subsubsection{Agent concepts}
In the field of artificial intelligence, an intelligent agent is often defined as an entity that perceives its environment and acts upon it autonomously in order to achieve its goals~\cite{Wooldridge1995IntelligentAgents}. Key characteristics of agents include autonomy, social ability, reactivity , and proactiveness~\cite{Wooldridge1995IntelligentAgents, singh2005service}. Agents encapsulate state about the world (beliefs or knowledge) and make decisions based on that state and their objectives. Unlike a service, which simply computes an output when invoked, an agent continually decides \emph{if}, \emph{when}, and \emph{what} actions to take next in order to meet its design objectives. This intrinsic goal-directed behavior means an agent can refuse or modify actions that conflict with its current goals or constraints, whereas a service will execute whenever called (assuming correct input). In a digital twin system, modeling a component as an agent means that component’s twin not only simulates the component’s state but also actively tries to optimize or control its performance towards certain goals (e.g., minimizing energy use, maximizing throughput) in coordination with other agents.

\subsubsection{Organizational structures}
Individual agents in a MAS do not operate in isolation, they are typically part of an organization of agents. MAS research has developed various organizational models to structure agent communities, ranging from flat peer-to-peer networks to hierarchical arrangements. Agents can dynamically assume roles (such as leader, coordinator, or participant) and establish relationships that facilitate coordinated action. For example, a group of agents may elect one of their number as a temporary coordinator for a task, or certain agents may act as brokers that match requests to providers. These organizational patterns can be fluid, changing as the system evolves, which gives MAS flexibility in how agents collaborate. A particularly relevant organizational concept for complex systems is the holarchy. In a holonic MAS, agents are arranged in a nested whole–part hierarchy of holons, where each agent is both an autonomous individual and part of a higher-level aggregate (holon). For instance, in a manufacturing context, each machine’s agent might be grouped into an agent representing a production cell, and multiple cell agents form part of a higher-level factory agent. This recursive hierarchy (holarchy) allows multi-level control and abstraction, combining the robustness of decentralized control with some advantages of top-down structure~\cite{HorlingLesser2005OrganizationalSurvey}. Organizational structuring helps align the MAS with the physical system’s structure; indeed, studies suggest using holarchies or team organizations to mirror the structure of coupled twin components~\cite{Human2023, Gil2024}. By organizing agents into appropriate structures (teams, federations, or holons), a MAS-based digital twin can manage complexity more effectively, ensuring that cooperation and communication occur at the right levels of granularity.

\subsubsection{Coordination and communication}
To function as a coherent system, agents must communicate and coordinate their actions. Unlike service invocation (which is typically a simple request/response), agent communication uses message-passing protocols that carry semantic intent. Agents send each other discrete messages (often asynchronously), using an Agent Communication Language (ACL) that defines performative verbs such as request, inform, propose, etc., to indicate the purpose of each message. Communication is often more dialog-oriented than in service calls: an agent may engage in a sequence of exchanges to negotiate or cooperate on a task. For example, a common interaction pattern is the Contract Net Protocol, where one agent (the manager) issues a call for proposals to others, who then bid with proposals; the manager evaluates bids and awards a contract to one or more agents to carry out the task. Such multi-step protocols allow decentralized task allocation and collaboration.

Coordination in MAS is achieved through these interactions rather than a central scheduler: agents coordinate by dynamically negotiating, cooperating, or even competing according to designed protocols and their own local decision logic. They can form coalitions or teams to achieve complex goals, and adjust their coordination patterns at runtime. This stands in contrast to the orchestration in service-oriented systems, which is typically predefined and static. The ability for agents to self-coordinate means that coherent global behavior can emerge from many local interactions. Notably, complex collective behaviors can arise without any single agent directing the whole process: a phenomenon known as emergent behavior. Classic examples include swarm or ant-colony systems where numerous simple agents collectively find near-optimal solutions to problems, as observed by \citeauthor{nwana1996software} in early multi-agent research~\cite{nwana1996software}. In a digital twin scenario, this could mean that multiple agents (each handling a part of the system) achieve a globally optimal schedule or adapt to disturbances through iterative local interactions, rather than following a fixed centrally-coded plan. Effective coordination mechanisms (negotiation strategies, voting schemes, mutual adjustment, etc.) are therefore critical MAS components to ensure that agent-based twins remain aligned with overall system objectives and do not work at cross purposes.

\subsubsection{Agent types and taxonomy}
Agents can be designed with different internal architectures and reasoning strategies, and the MAS literature classifies agents into various types based on these characteristics. Broadly, one major distinction is between reactive agents and deliberative agents. A reactive agent operates on a simple stimulus-response basis: it perceives an event or change and responds with a predefined or learned action, without maintaining an internal symbolic model of the environment. Reactive agents (for example, those based on Brooks’ subsumption architecture in robotics~\cite{Brooks1986Subsumption}) are fast and robust to dynamic changes, but they lack strategic planning ability and cannot easily pursue long-term goals beyond immediate reactions. In contrast, a deliberative agent (also called an intentional agent) maintains an explicit internal model and uses symbolic reasoning or planning to decide its actions. As \citeauthor{Wooldridge1995IntelligentAgents} defines, a deliberative agent is “one that contains an explicitly represented, symbolic model of the world, and in which decisions are made via logical (or at least pseudo-logical) reasoning”~\cite{Wooldridge1995IntelligentAgents}. Deliberative agents follow a sense–plan–act cycle and can consider future consequences, enabling them to handle complex, long-term tasks but sometimes at the cost of speed and responsiveness.

Within the deliberative category, a prominent architecture is the Belief–Desire–Intention (BDI) agent model. BDI agents maintain a structured mental state: \emph{beliefs} about the environment, \emph{desires} (goals or situations to achieve), and \emph{intentions} (the plans or courses of action to which the agent has committed)~\cite{RaoGeorgeff1995BDI}. The BDI framework provides a practical way to model decision-making: the agent continually updates its beliefs from percepts, generates options to achieve its desires, and selects which desires to commit to as intentions based on priority and feasibility. Intentions give the agent persistence, so it doesn’t constantly change plans at the slightest event, yet the agent can still react if something important changes (by updating beliefs and potentially revising its intentions)~\cite{RaoGeorgeff1995BDI}. BDI agents are well-suited for complex domains requiring reasoning about goals and plans, and have influenced many practical agent programming frameworks that implement this logic~\cite{RaoGeorgeff1995BDI}. A challenge with BDI and other deliberative agents is ensuring timely responses in dynamic environments, since deliberation can be computationally intensive.

Another important category is hybrid agents, which combine reactive and deliberative elements. In practice, purely reactive or purely deliberative approaches represent two extremes, and many real-world agent systems use a hybrid architecture to leverage the advantages of both~\cite{Wooldridge1995IntelligentAgents}. A common design is layering: a hybrid agent may include low-level reactive layers for immediate response to time-critical events, and higher-level deliberative layers for planning and goal management. For example, an autonomous vehicle’s agent might have a reactive layer for obstacle avoidance (reflexive, fast behavior) and a deliberative layer for route planning (strategic, goal-oriented behavior). The layers operate concurrently and inform each other~\cite{Muller1996InteRRaP, Ferguson1996TouringMachines}. Hybrid agents aim to achieve both the reactivity needed for dynamic environments and the foresight needed for complex tasks~\cite{palanca2023flexible, gandon2003combining}. The main difficulty is designing the coordination between layers (e.g., how the deliberative layer can override or guide the reactive layer without causing conflict).

Agents can be further specialized: for instance, holonic agents (in a holonic MAS) are those structured in the hierarchical manner described earlier, acting as both wholes and parts in the system. There are also learning agents that incorporate machine learning to improve their behavior over time, cognitive agents that attempt to emulate human-like reasoning, and semantic agents that leverage knowledge representation and ontologies. Figure~\ref{fig:agent-taxonomy} illustrates a taxonomy of intelligent agent types. In this diagram, “Intelligent Agents” are categorized into reactive, deliberative (with subtypes such as BDI, goal-based, and utility-based agents), hybrid, learning, cognitive, semantic, and organizational agents (with holonic agents as a subtype of the organizational category). This taxonomy highlights the diversity of agent designs available. For the purposes of digital twin architectures, the choice of agent type depends on the complexity of the task and environment: simple monitoring functions might be handled by reactive agents, whereas higher-level supervisory or optimization tasks might require deliberative or hybrid agents. Overall, MAS technology provides a rich toolkit of agent types and organizational patterns that can be tailored to create a modular, extensible digital twin system capable of addressing the limitations of monolithic designs.

\begin{figure}[ht]\centering
\begin{forest}
for tree={draw, rounded corners, align=center, inner sep=3pt, font=\sffamily}
[Intelligent Agents
[Reactive Agents]
[Deliberative Agents
[BDI Agents]
[Goal-Based Agents]
[Utility-Based Agents]
]
[Hybrid Agents]
% [Learning Agents]
% [Cognitive Agents]
% [Semantic Agents]
[Organizational Agents
[Holonic Agents]
]
]
\end{forest}
\caption{Taxonomy of agent types in multi-agent systems. Reactive, deliberative (including BDI, goal-based, utility-based), hybrid, and organizational (holonic) agents are key categories, alongside learning, cognitive, and semantic agents.}
\label{fig:agent-taxonomy}
\end{figure}

\subsection{Agent-based digital twin architectures}
Researchers have increasingly explored multi-agent system architectures as a basis for digital twins, especially in domains like smart manufacturing and cyber-physical systems. In these approaches, each agent typically represents a component or aspect of the physical system, and the digital twin emerges from the interactions of these agents. This paradigm tackles distributed intelligence and modularity. Agents encapsulate local decision-making and, in some cases, simulation models, cooperating to mirror real-world interactions. For example, in a factory setting, an agent-based twin can assign one agent per machine or cell, with higher-level coordinator agents managing workflows, which can improve production scheduling and responsiveness~\cite{latsou2021dt_mas_cpms}. Recent surveys note that MAS integration in digital twins can support adaptability, decentralized control, and real-time responsiveness across domains~\cite{Kalyani2024Survey}. Early work by \citeauthor{zhao2004agent} used agents to orchestrate interactive simulation components, separating simulation logic from control logic and illustrating how agents can support reusable components and complex scenario management~\cite{zhao2004agent}.

Despite these benefits, existing MAS-based digital twin solutions show gaps that limit fully modular and adaptive integration. Many implementations are domain-specific prototypes rather than general frameworks. They solve the particular problem at hand but are difficult to generalize or reconfigure. There is also a tendency to simplify coordination through fixed hierarchies or static team structures defined at design time. This limits adaptability when the real system changes or scales. Moreover, while agents can encapsulate models, integrating high-fidelity physics models or legacy simulators into an agent framework is non-trivial. Many agent-based twins emphasize control and decision-making more than continuous numeric simulation, which can limit hybrid use cases that need both.

An MAS-based digital twin architecture intended to be domain-agnostic and adaptive, combining autonomy and coordination with explicit support for modular integration of heterogeneous components is proposed. By incorporating a formal organizational structure and separating agent reasoning from simulation interfacing, the approach targets limitations in prior MAS-for-digital twin systems and aims to provide a more modular and extensible platform.

\subsection{Organizational models and coordination in MAS}
In multi-agent systems research, organizational models and coordination mechanisms address complexity in agent societies through concepts such as roles, groups, hierarchies, and coordination protocols. These abstractions are relevant for digital twin architectures with many agents because they impose structure on responsibilities and communication pathways. For example, the Agent-Group-Role (AGR) model formalizes an organization where each agent enacts roles within groups, supporting role-based communication and reusable interaction patterns \cite{ferber2003organizational_view}. Hierarchical organization paradigms arrange agents in layered control structures \cite{horling2004survey_org_paradigms}, which can mirror real-world organizational structures and support large-scale systems. Architecture-centric methods and middleware services can also help ensure properties like flexibility and reliability as the number of agents grows \cite{weyns2010architecture}.

At the same time, many MAS organizational frameworks are primarily design-time objects. Roles, groups, and norms are usually defined in advance. While some frameworks allow dynamic changes, they still assume a known organizational schema. In adaptive digital twins, the necessary roles or hierarchies may not be fully known prior, and the system may need to form new groups or coordination patterns at runtime. Existing MAS coordination techniques often focus on internal cooperation without explicitly addressing coupling to external simulators or evolving data streams. As a result, organizational models provide a foundation for structured coordination, but they do not ensure adaptivity and integration in digital twin settings.

The proposed architecture targets this gap by using organizational structure as a primary element while allowing runtime modification driven by the agents. This supports coordination that remains systematic through roles and protocols, while still enabling reorganization such as forming new groups to address emergent conditions.
